%% ============================================================================
\section{Lecture 11: Tsunami Amplification and Elasticity of Solids}
\label{sec:lec11}
%% ============================================================================
\date{25/06/2026}

\begin{center}
    \textit{Lecture date: 25/06/2026}
\end{center}

This lecture closes the discussion of water waves by asking what happens to a
tsunami as the ocean becomes shallow near shore. The same energy-flux logic is
then used as a bridge to solid-state estimates: atomic sizes, mass densities,
chemical bond energies, elastic moduli, sound speeds in solids, and the origin
of thermal expansion from anharmonic bond potentials.

%%----------------------------------------------------------------------------
\subsection{Tsunami Shoaling from Energy Flux Conservation}
\label{sec:lec11:tsunami-shoaling}
%%----------------------------------------------------------------------------

\subsubsection*{Physical Motivation}
Lecture~10 showed that a tsunami is a shallow-water gravity wave with phase
speed
\begin{equation}
    v_{\rm ph}=\sqrt{gH},
    \label{eq:lec11:shallow-speed}
\end{equation}
where $g$ is the gravitational acceleration and $H$ is the water depth. In the
open ocean $H$ is of order kilometers, while near shore it can be meters or
less. The wave therefore slows down as it approaches land. The question is:
what happens to its amplitude?

The useful conserved quantity is not the local energy density. As the depth
changes, the local wavelength and velocity change. What is conserved, in the
absence of strong dissipation or reflection, is the \emph{energy flux}: the
energy crossing a vertical line per unit time.

\subsubsection*{Wave Energy Density}
Let $Y$ denote the wave amplitude, i.e. the vertical displacement of the water
surface from its mean level. A common mistake is to write an energy proportional
to $Y$, but a wave energy must be quadratic in the amplitude. The gravitational
potential energy per unit horizontal area is estimated by lifting a column of
water of mass per area $\rho Y$ through a height of order $Y$:
\begin{equation}
    \frac{E}{A}\sim \rho gY^2 .
    \label{eq:lec11:water-energy-area}
\end{equation}
The missing numerical coefficient is order unity and is irrelevant for the
scaling argument. The important point is the quadratic dependence.

\subsubsection*{Flux Conservation}
The energy flux per unit crest length is the energy per unit area times the
speed:
\begin{equation}
    F_E \sim \rho gY^2 v_{\rm ph}.
    \label{eq:lec11:energy-flux}
\end{equation}
For a shallow-water wave, $v_{\rm ph}=\sqrt{gH}$, hence
\begin{equation}
    F_E \sim \rho gY^2\sqrt{gH}.
    \label{eq:lec11:energy-flux-shallow}
\end{equation}
If $\rho$ and $g$ are fixed and the flux is conserved, then
\begin{equation}
    Y^2H^{1/2}=\text{constant}.
    \label{eq:lec11:y2hhalf}
\end{equation}
Taking the square root gives
\begin{equation}
    \boxed{
    YH^{1/4}=\text{constant}.
    }
    \label{eq:lec11:green-law}
\end{equation}
\textit{Significance:} As depth decreases, the tsunami amplitude grows only as
$H^{-1/4}$ in the linear shallow-water regime; this weak power still becomes
large when ocean depth changes by orders of magnitude.

Comparing a deep-ocean region of depth $H_d$ and amplitude $Y_d$ with a shallower
region of depth $H_s$ and amplitude $Y_s$,
\begin{equation}
    \boxed{
    Y_s=Y_d\left(\frac{H_d}{H_s}\right)^{1/4}.
    }
    \label{eq:lec11:shoaling-ratio}
\end{equation}
\textit{Significance:} A one-metre tsunami in the deep ocean can become several
metres high near shore without requiring any new energy source.

\subsubsection*{Run-Up and the Nonlinear Shore Limit}
The linear estimate cannot be continued all the way to $H\to0$. When the wave
height becomes comparable to the water depth, the wave is no longer a small
perturbation. The relevant depth scale is then set by the wave height itself,
$H_s\sim Y_s$. Substituting this into flux conservation,
\begin{equation}
    Y_d^2H_d^{1/2}\sim Y_s^2Y_s^{1/2}=Y_s^{5/2}.
    \label{eq:lec11:runup-step}
\end{equation}
Raise both sides to the power $2/5$:
\begin{equation}
    \boxed{
    Y_s\sim Y_d^{4/5}H_d^{1/5}.
    }
    \label{eq:lec11:runup}
\end{equation}
\textit{Significance:} The final coastal run-up is a nonlinear estimate: once
the wave is as tall as the local water column, the small-amplitude shallow-water
description has broken down.

\nt{The distinction between shoaling and run-up matters. Equation
\eqref{eq:lec11:shoaling-ratio} describes linear amplification while the depth is
still externally prescribed. Equation~\eqref{eq:lec11:runup} estimates the
shoreline limit where the wave height itself sets the effective water depth.}

\subsubsection*{Wave Breaking and Shock Formation}
In a nonlinear shallow-water wave the local propagation speed depends on the
local depth. The crest sits on a larger effective water column than the trough,
so the crest can move faster. If the approach to shallow water is gradual enough,
the crest has time to overtake the lower part of the wave, and the wave breaks.

The same qualitative mechanism appears in other media. In the solar corona,
convection excites acoustic and magnetohydrodynamic waves that travel upward
into lower-density gas. Conservation of wave energy flux makes their amplitudes
grow. Once the perturbation is nonlinear, hotter or more compressed parts of the
wave propagate faster than cooler or less compressed parts, producing shocks.
Shock dissipation converts ordered wave energy into heat.

%%----------------------------------------------------------------------------
\subsection{Atomic Sizes and Solid Densities}
\label{sec:lec11:atomic-density}
%%----------------------------------------------------------------------------

\subsubsection*{Conceptual Background}
To estimate solid properties, we start with the simplest possible microscopic
picture: a solid is made of atoms packed a few Angstrom apart. The Bohr radius
of hydrogen is
\begin{equation}
    a_0\simeq0.53\,\text{\AA}.
    \label{eq:lec11:bohr}
\end{equation}
Most atomic radii from helium through uranium are of order
\begin{equation}
    a\simeq1.5\,\text{\AA}.
    \label{eq:lec11:atomic-radius}
\end{equation}
Thus an atom occupies a volume of order
\begin{equation}
    V_{\rm atom}\sim (3\,\text{\AA})^3.
    \label{eq:lec11:atomic-volume}
\end{equation}

\subsubsection*{Density Estimate}
Let $A$ be the atomic mass number. The mass per atom is $Am_p$, where $m_p$ is
the proton mass. The density is therefore
\begin{equation}
    \rho\sim \frac{Am_p}{(3\,\text{\AA})^3}.
    \label{eq:lec11:rho-step}
\end{equation}
Using $m_p=1.67\times10^{-24}\,\text{g}$ and
$3\,\text{\AA}=3\times10^{-8}\,\text{cm}$,
\begin{align}
    \rho
    &\sim
    \frac{A(1.67\times10^{-24}\,\text{g})}
         {(3\times10^{-8}\,\text{cm})^3}
    \notag\\
    &=
    \frac{1.67}{27}A\,\text{g\,cm}^{-3}
    \sim 0.06A\,\text{g\,cm}^{-3}.
    \label{eq:lec11:rho-number}
\end{align}
For heavy elements with $A\sim 100$--$200$, this gives densities of order
$10\,\text{g\,cm}^{-3}$, which is the right scale for metals. For light solids
the estimate gives values of order $1\,\text{g\,cm}^{-3}$.

\clm[clm:lec11:solid-density]{Atomic Packing Density Estimate}{
A solid made of atoms of radius $\sim1.5\,\text{\AA}$ has density
\begin{equation}
    \boxed{
    \rho\sim \frac{Am_p}{(3\,\text{\AA})^3}
    \sim 0.06A\,\text{g\,cm}^{-3}.
    }
    \label{eq:lec11:solid-density}
\end{equation}
}

\textit{Significance:} The observed density scale of ordinary matter follows
mainly from atomic size and nuclear mass; crystal structure changes only
order-unity packing factors.

\nt{Different lattices such as simple cubic, body-centred cubic, face-centred
cubic, hexagonal close packing, glasses, and quasicrystals change the packing
details, but they do not change the basic Angstrom-scale volume per atom.}

%%----------------------------------------------------------------------------
\subsection{Chemical Bond Energies}
\label{sec:lec11:bond-energies}
%%----------------------------------------------------------------------------

\subsubsection*{Physical Motivation}
Elasticity is controlled by how expensive it is to stretch or compress the
bonds between atoms or molecules. The first task is therefore to estimate the
energy scale of those bonds.

\subsubsection*{Ionic Bonds from Hydrogen Scaling}
The hydrogen ground-state energy is the Rydberg scale,
\begin{equation}
    E_{\rm H}=13.6\,\text{eV}.
    \label{eq:lec11:rydberg}
\end{equation}
An ionic bond, for example in NaCl, is roughly a Coulomb attraction between
oppositely charged ions separated by a distance of order a few Angstrom. Since
the Coulomb energy scales as $1/r$, compare the ionic separation
$r_{\rm ion}\sim3\,\text{\AA}$ to the Bohr radius $a_0\simeq0.5\,\text{\AA}$:
\begin{equation}
    E_{\rm ion}\sim E_{\rm H}\frac{a_0}{r_{\rm ion}}
    \sim 13.6\,\text{eV}\cdot\frac{0.5}{3}
    \sim 2\,\text{eV}.
    \label{eq:lec11:ionic-step}
\end{equation}
Depending on charges, geometry, and screening, real ionic bond energies are
typically a few electron-volts.

\begin{equation}
    \boxed{
    E_{\rm bond}\sim 1\text{--}10\,\text{eV}
    }
    \label{eq:lec11:bond-scale}
\end{equation}
\textit{Significance:} Ordinary chemical bonds are tied to the atomic Coulomb
scale, reduced from $13.6\,\text{eV}$ by the larger interatomic separation and
by screening or charge sharing.

\subsubsection*{Covalent and van der Waals Bonds}
A covalent bond is quantum mechanical: when two atoms are brought together,
the degenerate atomic energy levels split into a lower bonding state and a
higher antibonding state. Electrons occupy the lower state, so the total energy
is reduced. The splitting is again set by the overlap of Angstrom-scale
electronic wavefunctions, so the energy is also in the electron-volt range.

Van der Waals forces are weaker. Permanent dipole-dipole interactions scale as
$r^{-3}$ in energy for fixed dipole orientations. Induced-dipole interactions
between neutral nonpolar molecules are weaker still and scale as $r^{-6}$ in
the attractive part. Phenomenological Lennard-Jones potentials often combine
this attraction with a steep short-distance repulsion:
\begin{equation}
    U(r)=4\epsilon\left[\left(\frac{\sigma}{r}\right)^{12}
    -\left(\frac{\sigma}{r}\right)^6\right].
    \label{eq:lec11:lennard-jones}
\end{equation}

\nt{For mechanical properties, the relevant bond is the weakest bond that must
be distorted by the deformation. Ice, for example, has strong covalent O--H
bonds inside each water molecule, but much softer hydrogen bonds between
molecules. The elastic modulus is controlled mainly by the softer link.}

%%----------------------------------------------------------------------------
\subsection{Elastic Moduli}
\label{sec:lec11:elastic-moduli}
%%----------------------------------------------------------------------------

\subsubsection*{Mathematical Setup}
Elastic moduli compare stress to strain. Stress has units of force per area,
equivalently energy per volume. Strain is dimensionless. Therefore every elastic
modulus has units of energy density.

\dfn[dfn:lec11:young]{Young Modulus}{
For a rod of length $L$ and cross-sectional area $A$, pulled by a force $F$,
the tensile stress is $F/A$ and the longitudinal strain is $\Delta L/L$. The
Young modulus $Y_{\rm Young}$ is defined by
\begin{equation}
    \boxed{
    \frac{F}{A}=Y_{\rm Young}\frac{\Delta L}{L}.
    }
    \label{eq:lec11:young-def}
\end{equation}
}

\textit{Significance:} $Y_{\rm Young}$ measures how much stress is required to
produce a given fractional extension.

\dfn[dfn:lec11:bulk-shear]{Bulk and Shear Moduli}{
The bulk modulus $K$ describes compression:
\begin{equation}
    \boxed{
    \Delta P=-K\frac{\Delta V}{V}.
    }
    \label{eq:lec11:bulk-def}
\end{equation}
The shear modulus $G_{\rm sh}$ describes shape change at nearly fixed volume:
\begin{equation}
    \boxed{
    \frac{F}{A}=G_{\rm sh}\frac{\Delta x}{L}.
    }
    \label{eq:lec11:shear-def}
\end{equation}
}

\textit{Significance:} Bulk response, shear response, and tensile response are
different experiments, but their magnitudes are all set by microscopic bond
curvatures divided by atomic volumes.

\subsubsection*{Order-of-Magnitude Modulus from Bond Energy}
Let $E_b$ be a typical bond energy and let $a$ be the interatomic spacing.
Since modulus has dimensions of energy density,
\begin{equation}
    Y_{\rm Young}\sim K\sim G_{\rm sh}\sim \frac{E_b}{a^3}.
    \label{eq:lec11:modulus-energy-density-step}
\end{equation}
With $E_b\sim4\,\text{eV}$ and $a\sim3\,\text{\AA}$,
\begin{align}
    \frac{E_b}{a^3}
    &\sim
    \frac{4(1.6\times10^{-12}\,\text{erg})}
         {(3\times10^{-8}\,\text{cm})^3}
    \notag\\
    &\sim
    2\times10^{11}\,\text{erg\,cm}^{-3}
    =
    2\times10^{10}\,\text{Pa}.
    \label{eq:lec11:modulus-number}
\end{align}
Real elastic moduli range around this value: ice is softer, typical metals are
$10^{10}$--$10^{11}\,\text{Pa}$, diamond is much stiffer, and graphite is very
anisotropic because strong in-plane covalent bonds coexist with weak interlayer
bonds.

\begin{equation}
    \boxed{
    M_{\rm elastic}\sim \frac{E_b}{a^3}
    \sim 10^{10}\text{--}10^{11}\,\text{Pa}.
    }
    \label{eq:lec11:elastic-scale}
\end{equation}
\textit{Significance:} The stiffness of ordinary solids is the chemical bond
energy packed into an atomic volume.

%%----------------------------------------------------------------------------
\subsection{Elastic Modulus from the Bond Potential}
\label{sec:lec11:bond-curvature}
%%----------------------------------------------------------------------------

\subsubsection*{Mathematical Setup}
The energy-density estimate explains the scale. To see where the modulus comes
from microscopically, expand the bond potential near its equilibrium separation.
Let $x$ be the bond displacement from equilibrium and let $d$ be the equilibrium
bond length. Since equilibrium is a minimum, the linear term vanishes:
\begin{equation}
    E(x)=E_0+\frac{1}{2}kx^2+\cdots,
    \qquad
    k=\left.\dv[2]{E}{x}\right|_{0}.
    \label{eq:lec11:harmonic-bond}
\end{equation}
The force in one bond is
\begin{equation}
    F_1=-\dv{E}{x}=-kx.
    \label{eq:lec11:bond-force}
\end{equation}

Consider a rod stretched by $\Delta L$. If all microscopic bonds share the
strain, then the displacement of one bond is
\begin{equation}
    x=\Delta d=d\frac{\Delta L}{L}.
    \label{eq:lec11:micro-strain}
\end{equation}
The number of parallel bonds crossing a section of area $A$ is of order
$A/a^2$, so the stress is
\begin{align}
    \frac{F}{A}
    &\sim
    \frac{1}{A}\frac{A}{a^2}k\Delta d
    =
    \frac{k}{a^2}d\frac{\Delta L}{L}.
    \label{eq:lec11:stress-micro}
\end{align}
Comparing with Eq.~\eqref{eq:lec11:young-def},
\begin{equation}
    \boxed{
    Y_{\rm Young}\sim \frac{k d}{a^2}
    =
    \frac{d}{a^2}\left.\dv[2]{E}{x}\right|_0
    \sim \frac{E_b}{a^3}.
    }
    \label{eq:lec11:young-curvature}
\end{equation}
\textit{Significance:} The elastic modulus is the curvature of the interatomic
potential, multiplied by the geometry of how many bonds transmit force across a
unit area.

In the last step we used $d\sim a$ and estimated the potential curvature as
$k\sim E_b/a^2$, because the energy changes by order $E_b$ over an atomic
distance.

%%----------------------------------------------------------------------------
\subsection{Sound Speed in a Solid}
\label{sec:lec11:sound-solid}
%%----------------------------------------------------------------------------

\subsubsection*{Physical Motivation}
Sound speed in a gas was controlled by $\sqrt{P/\rho}$ because pressure was the
restoring stress. In a solid, pressure is not the relevant microscopic stiffness.
The restoring stress is elastic, so the pressure scale is replaced by an
elastic modulus.

Dimensional analysis gives
\begin{equation}
    c_s\sim\sqrt{\frac{M_{\rm elastic}}{\rho}}.
    \label{eq:lec11:sound-solid-step}
\end{equation}
Using $M_{\rm elastic}\sim E_b/a^3$ and
$\rho\sim m_{\rm atom}/a^3$,
\begin{equation}
    \boxed{
    c_s\sim
    \sqrt{\frac{E_b/a^3}{m_{\rm atom}/a^3}}
    =
    \sqrt{\frac{E_b}{m_{\rm atom}}}.
    }
    \label{eq:lec11:sound-solid}
\end{equation}
\textit{Significance:} The sound speed in a solid is the characteristic atomic
vibration speed set by bond energy per atom divided by atomic mass.

\ex[ex:lec11:sound-solid]{Worked Example: Solid Sound Speed}{
Take $E_b\sim4\,\text{eV}=6.4\times10^{-19}\,\text{J}$ and an atomic mass
$m_{\rm atom}\sim30m_p\sim5\times10^{-26}\,\text{kg}$. Then
\[
    c_s\sim\sqrt{\frac{6.4\times10^{-19}}{5\times10^{-26}}}
    \sim \sqrt{1.3\times10^7}
    \sim 4\times10^3\,\text{m/s}.
    \]
This is the correct kilometre-per-second scale for sound in many solids.
}

%%----------------------------------------------------------------------------
\subsection{Thermal Expansion and Anharmonicity}
\label{sec:lec11:thermal-expansion}
%%----------------------------------------------------------------------------

\subsubsection*{Conceptual Background}
Thermal expansion is not caused merely by atoms moving more when they are hot.
In a perfectly harmonic potential,
\begin{equation}
    U(x)=\frac{1}{2}kx^2,
    \label{eq:lec11:harmonic-potential}
\end{equation}
the Boltzmann distribution is symmetric:
\begin{equation}
    P(x)\propto \exp\left[-\frac{kx^2}{2k_BT}\right].
    \label{eq:lec11:boltzmann-harmonic}
\end{equation}
The mean displacement is therefore
\begin{equation}
    \langle x\rangle =0
    \label{eq:lec11:harmonic-mean}
\end{equation}
at every temperature, even though $\langle x^2\rangle$ grows with $T$.

\subsubsection*{Why Anharmonicity is Required}
To make the average bond length increase with temperature, the potential must
be asymmetric. Expand the bond potential to the next order:
\begin{equation}
    U(x)=E_0+\frac{1}{2}kx^2+\frac{1}{6}\lambda x^3+\cdots,
    \qquad
    \lambda=\left.\dv[3]{U}{x}\right|_0 .
    \label{eq:lec11:anharmonic}
\end{equation}
The cubic term is the first term that can distinguish positive and negative
displacements. Its natural size is estimated as
\begin{equation}
    \lambda\sim \frac{E_b}{a^3},
    \label{eq:lec11:cubic-scale}
\end{equation}
while the harmonic curvature is
\begin{equation}
    k\sim\frac{E_b}{a^2}.
    \label{eq:lec11:k-scale}
\end{equation}

\clm[clm:lec11:thermal-expansion]{Origin of Thermal Expansion}{
A purely harmonic crystal does not thermally expand. Thermal expansion requires
anharmonicity in the interatomic potential:
\begin{equation}
    \boxed{
    \alpha_L\equiv\frac{1}{L}\dv{L}{T}
    \sim \frac{k_B}{E_b}.
    }
    \label{eq:lec11:alpha-thermal}
\end{equation}
}

\textit{Significance:} The fractional expansion per kelvin is small because
room-temperature thermal energy $k_BT$ is much smaller than a chemical bond
energy.

For $E_b\sim1$--$5\,\text{eV}$,
\begin{equation}
    \alpha_L\sim\frac{8.6\times10^{-5}\,\text{eV/K}}{1\text{--}5\,\text{eV}}
    \sim 10^{-5}\,\text{K}^{-1},
    \label{eq:lec11:alpha-number}
\end{equation}
which is the observed scale for ordinary solids.

\nt{The lecture stopped just before doing the full thermal-expansion estimate.
The key setup is already fixed: the second derivative of the bond potential
sets elastic moduli, while the third derivative sets the leading anharmonic
correction responsible for expansion.}

%%----------------------------------------------------------------------------
\subsection{Summary}
\label{sec:lec11:summary}
%%----------------------------------------------------------------------------

\begin{enumerate}
    \item Tsunami amplitude growth follows from energy flux conservation:
    $F_E\sim\rho gY^2\sqrt{gH}$, hence $YH^{1/4}=\text{constant}$ in the linear
    shallow-water regime.
    \item Near shore the small-amplitude description fails when $Y\sim H$,
    giving the nonlinear run-up estimate $Y_s\sim Y_d^{4/5}H_d^{1/5}$.
    \item Ordinary solid densities are set by Angstrom-scale atomic volumes:
    $\rho\sim Am_p/(3\,\text{\AA})^3$.
    \item Chemical bond energies are typically $1$--$10\,\text{eV}$, inherited
    from Coulomb physics at atomic length scales.
    \item Elastic moduli are bond energy densities:
    $M_{\rm elastic}\sim E_b/a^3\sim10^{10}$--$10^{11}\,\text{Pa}$.
    \item The sound speed in a solid is
    $c_s\sim\sqrt{M_{\rm elastic}/\rho}\sim\sqrt{E_b/m_{\rm atom}}$.
    \item Thermal expansion is an anharmonic effect: harmonic motion increases
    $\langle x^2\rangle$ but leaves $\langle x\rangle$ unchanged.
\end{enumerate}
